\begin{abstract}

Physical agents such as autonomous vehicles and robot manipulators are increasingly built by learning from data, yet they are deployed in an open world that keeps presenting conditions their training data do not contain. This thesis calls this challenge the \emph{curse of reality} and studies how to make physical agents robust to it along the perceive--reason--act loop.

For perception, the thesis proposes to describe driving scenes by 3D semantic occupancy instead of bounding boxes. It introduces OccNet, a camera-based network that reconstructs occupancy through a cascade of voxel decoders, and OpenOcc, a dense occupancy benchmark, and shows that occupancy describes irregular objects better than boxes, transfers to detection and segmentation, and lowers the collision rate of a planner. It further scales occupancy supervision with OpenScene, a large redistribution of real-world driving data that has served community challenges on end-to-end driving and world models.

For reasoning, the thesis introduces Graph Visual Question Answering, which links perception, prediction and planning questions through logical dependencies, together with the DriveLM datasets and a vision-language baseline. Graph-structured reasoning improves zero-shot generalisation to a new sensor configuration and to unseen objects, although it brings little benefit in familiar situations.

For acting, the thesis proposes Centaur, which lets an end-to-end driving planner correct itself at deployment by test-time training on Cluster Entropy, a label-free measure of the planner's uncertainty. Centaur raises the driving score on the NAVSIM benchmark from 90.3 to 92.6 without hand-designed rules, and a new safety-critical split exposes the remaining gap to human driving. For robot manipulation, the thesis presents $\chi_0$, which traces the fragility of real-world policies to inconsistencies between demonstrations, model and execution, and addresses them with Model Arithmetic, Stage Advantage and Train--Deploy Alignment. With about 20 hours of data per task and eight GPUs, $\chi_0$ surpasses the success rate of a state-of-the-art generalist policy by nearly 250\% on garment manipulation and runs for 24 hours without interruption.

Drawing on these contributions and on the author's other publications, the thesis distils a benchmark-driven research methodology. It decomposes the curse of reality into perceptual, semantic, behavioural and deployment shift, and proceeds through a cycle that turns failures observed in the field into public benchmarks, simple and scalable baselines, community challenges, and self-correction at deployment. Synthesising the findings, the thesis argues that a self-correcting physical agent needs an open and queryable representation, fast and slow decision paths, a label-free monitor of its own competence, adaptation at deployment, and a learning loop that feeds residual failures back into data and benchmarks.

\vspace{1em}
\noindent(An abstract of 416 words)

\end{abstract}
